\documentclass{article}

\usepackage{amsthm}
\usepackage{thmtools}
\usepackage[dvipsnames]{xcolor}
\usepackage{lipsum}
\usepackage{amsmath}

\newtheorem{theorem}{Theorem}[section]
\newtheorem{result}{Result}[theorem]
\newtheorem{lemma}[theorem]{Lemma}
\newtheorem*{plain}{Plain}

\theoremstyle{definition}
\newtheorem{definition}{Definition}[section]

\theoremstyle{remark}
\newtheorem{remark}{Remark}

\newtheoremstyle{myStyle}
				{1cm}{}
				{\itshape}{}
				{\bfseries}{:}
				{\newline}{}
\theoremstyle{myStyle}
\newtheorem{custom}{Custom}

\declaretheorem[shaded={bgcolor=SkyBlue},
				name=ShadedThm]{shTheorem}
				
\declaretheorem[shaded={rulecolor=YellowGreen,
						rulewidth=2pt,
						bgcolor=white}]{ruledThm}
						
%\declaretheorem[thmbox=L]{LBoxThm}
%\declaretheorem[thmbox=M]{MBoxThm}
%\declaretheorem[thmbox=S]{SBoxThm}

\renewcommand\qedsymbol{\fbox{Ended}}

\begin{document}
	{
		\noindent
		\Huge
		\textbf{Theorems and Proofs}	
	}

	\vspace{1cm}
	
	\section{First Section}
	\begin{theorem}[TheoremName]
		Put the content of the theorem here !!
	\end{theorem}
	
	\section{Second Section}	
	\begin{theorem}
		Put the content of the theorem here !!
	\end{theorem}
	\begin{theorem}
		Put the content of the theorem here !!
	\end{theorem}
	
	\begin{result}
		The new content of the result env. !!
	\end{result}
	
	\begin{lemma}
		Put the content of the theorem-like env. here.
	\end{lemma}
	
	\begin{theorem}
		Put the content of the theorem here !!
	\end{theorem}
	
	\begin{plain}
		This is an unnumbered theorem-like env.
	\end{plain}
	
	\section{Theorem Styles}
	\begin{definition}
		Put your definition here !!
	\end{definition}
	\begin{remark}
		Put your remark here !!
	\end{remark}
	
	\section{Custom Theorem Style}
	\begin{custom}
		Custom Theorem Style Environment.
	\end{custom}
	\begin{custom}
		Custom Theorem Style Environment.
	\end{custom}
	
	\pagebreak
	
	\section{thmtools Package}
	\subsection{shadethm Package Support}
	\begin{shTheorem}
		For every prime $p$, there then is a prime 
			$p'>p$. In Particular, there are infinitely 			man primes.		
	\end{shTheorem}
	\begin{ruledThm}
	For every prime $p$, there then is a prime 
			$p'>p$. In Particular, there are infinitely 			man primes.
	\end{ruledThm}
	
	\vspace{2cm}
	
%	\subsection{thmbox Package Support}
%	\begin{LBoxThm}
%		For every prime $p$, there then is a prime 
%			$p'>p$. In Particular, there are infinitely 			man primes.
%	\end{LBoxThm}
%	\begin{MBoxThm}
%		For every prime $p$, there then is a prime 
%			$p'>p$. In Particular, there are infinitely 			man primes.
%	\end{MBoxThm}
%	\begin{SBoxThm}
%		For every prime $p$, there then is a prime 
%			$p'>p$. In Particular, there are infinitely 			man primes.
%	\end{SBoxThm}
	
	\pagebreak
	
	\section{Proofs}
	The Proof Environment is Like:
	\begin{proof}[myNewProof]
		Put the content of the proof here !!
		\renewcommand{\qedsymbol}{}
	\end{proof}
	
	\lipsum[1] \qed
	
	\begin{proof}
		\renewcommand\qedsymbol{$\clubsuit$}
		\begin{equation*}
			F(x) = 2x + y
			\qedhere
		\end{equation*}
	\end{proof}
	
\end{document}